\section{Root-to-Node Path}
\label{sec:trees-root-to-node-path}

\problemstatement{Given the root of a binary tree and a target node value
(assume values are unique), return the sequence of node values on the path
from the root down to the target.}

\begin{patternbox}
\emph{``Path from the root to a node,'' ``ancestors of a node''} is a
backtracking DFS: descend adding the current node to a running path; if
this node is the target, stop and keep the path; if neither subtree
contains the target, \textbf{remove the current node} before returning. The
add-then-remove-on-failure rhythm is the essence of backtracking.
\end{patternbox}

\subsection*{Understanding the problem carefully}

There is exactly one path from the root to any node in a tree. We build it
incrementally: push the current node onto the path, then check ``is this
the target, or is the target somewhere below on the left or right?'' If
yes, this node is genuinely on the path and stays. If the target is not
here and not in either subtree, this node was a wrong turn and must be
popped so it does not pollute the path for other branches.

\begin{ideabox}[Add on Entry, Pop on a Dead End]
\textsc{Find}(node, target, path): if node is null return false; push
node.val onto path; if node.val equals target return true; if
\textsc{Find}(left) or \textsc{Find}(right) returns true, return true
(node stays on the path); otherwise pop node.val (dead end) and return
false.
\end{ideabox}

\subsection*{Algorithm}

\begin{enumerate}
  \item \textsc{Find}(\textit{node}, \textit{target}, \textit{path}):
  \begin{enumerate}
    \item If \textit{node} is null, return false.
    \item Push \textit{node}.val onto \textit{path}.
    \item If \textit{node}.val $=$ \textit{target}, return true.
    \item If \textsc{Find}(left) \textbf{or} \textsc{Find}(right),
          return true.
    \item Pop \textit{node}.val; return false.
  \end{enumerate}
\end{enumerate}

\subsection*{Dry run}

\dryrun{Tree: root $1$; left $2$ (children $4,5$); right $3$. Target $5$.}

\begin{itemize}
  \item Push $1$. Not $5$. Recurse left into $2$.
  \item Push $2$. Not $5$. Recurse left into $4$: push $4$, not $5$, no
        children, pop $4$, return false. Recurse right into $5$: push
        $5$, equals target, return true.
  \item $2$ keeps ($5$ found below); $1$ keeps.
  \item Path: $1,2,5$.
\end{itemize}

\subsection*{C++ implementation}

\begin{lstlisting}
#include <bits/stdc++.h>
using namespace std;

struct TreeNode {
    int val;
    TreeNode* left;
    TreeNode* right;
    TreeNode(int v) : val(v), left(nullptr), right(nullptr) {}
};

bool findPath(TreeNode* node, int target, vector<int>& path) {
    if (node == nullptr) return false;
    path.push_back(node->val);
    if (node->val == target) return true;
    if (findPath(node->left, target, path) ||
        findPath(node->right, target, path))
        return true;
    path.pop_back();                 // dead end: backtrack
    return false;
}

int main() {
    TreeNode* root = new TreeNode(1);
    root->left = new TreeNode(2);
    root->right = new TreeNode(3);
    root->left->left = new TreeNode(4);
    root->left->right = new TreeNode(5);

    vector<int> path;
    findPath(root, 5, path);
    for (int x : path) cout << x << " ";
    cout << "\n";
    return 0;
}
\end{lstlisting}

\begin{complexitybox}
\textbf{Time Complexity.} $O(n)$ -- in the worst case every node is
visited once. \textbf{Space Complexity.} $O(h)$ for the recursion stack
and the path.
\end{complexitybox}

\begin{takeawaybox}
The push-on-entry / pop-on-failure pattern is backtracking in its simplest
tree form, and it is the conceptual basis for the next section's lowest
common ancestor and for any ``list the ancestors'' query.
\end{takeawaybox}
